\subsection{Estado del arte discriminante y auditoría de novedad}
\label{subsec:tf-discriminant}

La Sección~\ref{subsec:tf-synthesis-gap} localizó la interfaz abierta entre composición
nominal, certificado mecánico, recuperación y tráfico. Esa formulación todavía es
cómoda, porque no se ha sometido a la prueba que decide si el problema focal existe:
comparados uno a uno, ¿qué cubre cada trabajo próximo y qué deja fuera? La
Figura~\ref{fig:lit-methodological-map} sitúa los 32 trabajos discriminantes del núcleo
canónico sobre los ejes de autoridad de decisión y explicitud del mecanismo. El patrón
que muestra es una división por familia: los trabajos de coalición y MRTA quedan del
lado de la decisión lógica y los de transporte cooperativo del lado de la ejecución
física. La ocupación de la franja intermedia no se interpreta aquí: el mapa codifica
32 de los 59 trabajos canónicos y no se comprobó si los 27 restantes la poblarían
(Anexo~\ref{app:review-full-detail}). El argumento de este apartado no descansa en esa
densidad, sino en la comparación trabajo a trabajo que sigue.

% Figura 1 del consolidado académico: mapa metodológico de los 32 trabajos
% discriminantes del núcleo canónico. El cuerpo TikZ y los macros \mappointL y
% \mappointR se portan desde `thesis-reviews/academic_literature_review (1).tex`
% (líneas 60-148), cuyo conjunto de trabajos coincide con el del PDF vigente.
\begin{figure}[!htbp]
  \centering
  \makebox[\textwidth][c]{\resizebox{0.98\textwidth}{!}{%
\begin{tikzpicture}[x=80mm,y=54mm,font=\sffamily]
  \fill[softblue] (-1,-1) rectangle (0,0);
  \fill[softgreen] (0,-1) rectangle (1,0);
  \fill[softplum!55!white] (-1,0) rectangle (0,1);
  \fill[softwarm] (0,0) rectangle (1,1);
  \foreach \x in {-1,-.75,-.5,-.25,0,.25,.5,.75,1} \draw[rule,line width=.28pt] (\x,-1)--(\x,1);
  \foreach \y in {-1,-.75,-.5,-.25,0,.25,.5,.75,1} \draw[rule,line width=.28pt] (-1,\y)--(1,\y);
  \draw[-{Stealth[length=2.3mm]},ink,line width=.68pt] (-1.04,0)--(1.075,0);
  \draw[-{Stealth[length=2.3mm]},ink,line width=.68pt] (0,-1.04)--(0,1.09);
  \node[anchor=south,font=\bfseries\fontsize{7.2}{7.5}\selectfont] at (0,1.10) {data-driven};
  \node[anchor=north,font=\bfseries\fontsize{7.2}{7.5}\selectfont] at (0,-1.09) {white-box};
  \node[rotate=90,anchor=south,font=\bfseries\fontsize{7}{7.3}\selectfont] at (-1.12,0) {descentralizado};
  \node[rotate=90,anchor=north,font=\bfseries\fontsize{7}{7.3}\selectfont] at (1.12,0) {centralizado};
  \node[align=center,font=\bfseries\fontsize{6.35}{6.65}\selectfont] at (-.60,.90) {distribuido\\[-.35mm]data-driven};
  \node[align=center,font=\bfseries\fontsize{6.35}{6.65}\selectfont] at (.60,.90) {centralizado\\[-.35mm]data-driven};
  \node[align=center,font=\bfseries\fontsize{6.35}{6.65}\selectfont] at (-.60,-.90) {distribuido\\[-.35mm]white-box};
  \node[align=center,font=\bfseries\fontsize{6.35}{6.65}\selectfont] at (.60,-.90) {centralizado\\[-.35mm]white-box};

  % data-driven, distributed / mixed
  \mappointL{-0.82}{0.72}{Bezerra}{2025}{B}{20}{learncol}
  \mappointL{-0.64}{0.55}{Huo}{2025}{H}{19}{learncol}
  \mappointL{-0.46}{0.73}{Dai}{2024}{D}{17}{learncol}
  \mappointL{-0.78}{0.35}{Shibata-RL}{2023}{S}{16}{learncol}
  \mappointL{-0.43}{0.32}{Pandit}{2026}{P}{28}{learncol}
  \mappointR{-0.13}{0.50}{Eoh/Park}{2021}{E}{29}{learncol}

  % data-driven, centralized/mixed
  \mappointL{0.34}{0.64}{SADCHER}{2026}{S}{25}{learncol}
  \mappointL{0.62}{0.42}{CF-HMRTA}{2025}{C}{21}{searchcol}

  % white-box distributed - staggered
  \mappointL{-0.90}{-0.18}{Pereira}{2004}{P}{1}{controlcol}
  \mappointL{-0.78}{-0.34}{Parker/Tang}{2006}{PT}{2}{foundcol}
  \mappointL{-0.88}{-0.54}{Vig/Adams}{2006}{V}{3}{gamecol}
  \mappointL{-0.69}{-0.17}{CBBA}{2009}{C}{4}{gamecol}
  \mappointL{-0.64}{-0.39}{Service}{2011}{S}{5}{foundcol}
  \mappointL{-0.57}{-0.60}{Chen/Sun}{2012}{CS}{6}{gamecol}
  \mappointL{-0.45}{-0.22}{Jang}{2018}{J}{8}{gamecol}
  \mappointL{-0.39}{-0.47}{Dutta}{2019}{D}{10}{gamecol}
  \mappointL{-0.36}{-0.67}{Dutta}{2021}{D}{12}{gamecol}
  \mappointL{-0.14}{-0.28}{Ebel}{2024}{E}{18}{controlcol}
  \mappointR{-0.03}{-0.50}{Shibata}{2023}{S}{14}{controlcol}
  \mappointR{-0.04}{-0.71}{Fukao}{2025}{F}{22}{controlcol}
  \mappointL{-0.48}{-0.90}{GRAPE-S}{2026}{G}{24}{gamecol}

  % white-box central/mixed
  \mappointL{0.08}{-0.18}{Alonso-Mora}{2017}{AM}{7}{controlcol}
  \mappointL{0.20}{-0.40}{Machado}{2019}{M}{9}{controlcol}
  \mappointL{0.31}{-0.65}{Verma repl.}{2019}{V}{11}{controlcol}
  \mappointL{0.47}{-0.22}{Elaamery}{2021}{E}{13}{controlcol}
  \mappointL{0.59}{-0.44}{Gong}{2023}{G}{15}{controlcol}
  \mappointL{0.76}{-0.22}{Rizzo}{2020}{R}{30}{controlcol}
  \mappointR{0.89}{-0.40}{Loh}{2012}{L}{31}{controlcol}
  \mappointR{0.83}{-0.64}{Yufka}{2015}{Y}{32}{controlcol}
  \mappointR{0.61}{-0.73}{Muhammed}{2026}{M}{26}{controlcol}
  \mappointR{0.34}{-0.88}{Sahu}{2026}{S}{27}{controlcol}
  \mappointL{0.02}{-0.90}{Tian}{2025}{T}{23}{controlcol}

  \node[star,star points=5,star point ratio=2.1,fill=orange,draw=white,line width=.8pt,minimum size=8.5mm] at (-0.18,-0.78) {};
  \node[anchor=west,font=\bfseries\fontsize{7.0}{7.4}\selectfont,align=left] at (-0.09,-0.75) {Este TFM\\integración};
\end{tikzpicture}%
  }}
  \par\vspace{1mm}
  {\centering\fontsize{8.35}{8.9}\selectfont
  \chip{gamecol}{juegos/mercados}\hspace{5.5mm}
  \chip{controlcol}{control/consenso}\hspace{5.5mm}
  \chip{searchcol}{búsqueda/heurística}\hspace{5.5mm}
  \chip{learncol}{política aprendida}\hspace{5.5mm}
  \chip{foundcol}{fundacional}\par}
  \caption[Mapa metodológico del núcleo canónico.]{Mapa metodológico de 32
  trabajos discriminantes del núcleo canónico. El eje horizontal codifica
  autoridad de decisión (descentralizado $\rightarrow$ centralizado) y el
  vertical la explicitud del mecanismo (white-box $\rightarrow$ data-driven); el
  color identifica la familia primaria. Los trabajos de coalición y MRTA se
  concentran en la decisión lógica y los de transporte cooperativo en la
  ejecución física. Las posiciones son ordinales y se desplazan solo para evitar
  solapes: no representan calidad ni rendimiento.}
  \label{fig:lit-methodological-map}
  \viusource{elaboración propia a partir de la codificación del núcleo canónico}
\end{figure}


El primer frente resuelve la coalición lógica mediante emparejamientos uno a muchos,
subastas y juegos hedónicos
\parencite{choiBrunetHow2009CBBA,serviceAdams2011Coalition,chenSun2012Coalition,jangShinTsourdos2018GRAPE,duttaAsaithambi2019Bipartite,dutta2021hedonic,diehlAdams2026Grapes};
el segundo resuelve la ejecución física con formulación, control y validación
experimental
\parencite{alonsomora2017transport,shibata2023event,ebel2024cooperative,fukao2025swarm,tian2025irregular,muhammed2026decentralizedMpc}.
La brecha no puede formularse como «falta control distribuido de carga», que sería
falso, sino como una interfaz que ninguno de los dos frentes atraviesa.

% Figura 2(a) y 2(b) del consolidado académico. La matriz se porta desde
% `thesis-reviews/academic_literature_review (1).tex` (líneas 176-217) y se le
% añade la columna C ---nivel de centralización C0--C4--- que incorpora el PDF
% vigente. Las modalidades físicas conservan los rótulos del original.
\begin{figure}[!htbp]
  \centering
  \makebox[\textwidth][c]{\resizebox{0.99\textwidth}{!}{%
\begin{minipage}{\linewidth}
\renewcommand{\arraystretch}{1.45}\setlength{\tabcolsep}{2.25pt}\fontsize{7.55}{8.0}\selectfont
\begin{tabularx}{\linewidth}{J{27mm}*{8}{K{10.75mm}}K{9.6mm}K{9.2mm}}
\toprule
\textbf{Trabajo} & \rotatebox{90}{\textbf{heterog.}} & \rotatebox{90}{\textbf{coal. var.}} & \rotatebox{90}{\textbf{reclut. local}} & \rotatebox{90}{\textbf{shared-load}} & \rotatebox{90}{\textbf{cert. física}} & \rotatebox{90}{\textbf{exec. dist.}} & \rotatebox{90}{\textbf{recovery}} & \rotatebox{90}{\textbf{replacement}} & \textbf{Ev.} & \textbf{C}\\
\midrule
Jang 2018 [8] & \capno & \capyes & \capyes & \capno & \capno & \capyes & \capno & \capno & T+S & C4\\
Dutta 2021 [12] & \capyes & \capyes & \capyes & \capno & \capno & \capyes & \capno & \capno & T+S & C4\\
GRAPE-S 2026 [24] & \capyes & \capyes & \cappart & \capno & \capno & \capyes & \capno & \capno & T+S & C3--4\\
CF-HMRTA 2025 [21] & \capyes & \capyes & \capno & \capno & \capno & \capno & \capno & \capno & T+S & C0--1\\
Huo 2025 [19] & \capyes & \capyes & \capyes & \capno & \capno & \capyes & \cappart & \capno & T+S & C3--4\\
Bezerra 2025 [20] & \capyes & \capyes & \capyes & \capno & \capno & \capyes & \cappart & \capno & S & C4\\
Alonso-Mora 2017 [7] & \capno & \cappart & \capno & \capyes & \capyes & \cappart & \cappart & \capno & T+S+P & C0--1\\
Ebel 2024 [18] & \capno & \capno & \capno & \capyes & \capyes & \capyes & \capno & \capno & T+S+P & C3\\
Shibata 2023 [14] & \capno & \cappart & \capno & \capyes & \capyes & \capyes & \cappart & \cappart & S+P & C4\\
Muhammed 2026 [26] & \capno & \capno & \capno & \capyes & \capyes & \capyes & \capno & \capno & T+S+P & C3\\
Verma repl. 2019 [11] & \capno & \cappart & \capno & \capyes & \cappart & \capno & \capyes & \capyes & S+P & C0--1\\
Sahu 2026 [27] & \capno & \cappart & \capno & \capyes & \cappart & \cappart & \capyes & \capyes & S+P & C1--2\\
\bottomrule
\end{tabularx}
\end{minipage}%
  }}
  \par\vspace{2mm}
  \makebox[\textwidth][c]{\resizebox{0.92\textwidth}{!}{%
\begin{tikzpicture}[x=1mm,y=1mm,font=\sffamily\fontsize{6.25}{6.65}\selectfont]
\tikzset{bx/.style={rounded corners=1pt,draw=rule,line width=.35pt,minimum width=41.5mm,minimum height=17mm,align=center}}
\node[bx,fill=softblue](s)at(21,10){\textbf{Supported cargo}\\carga soportada\\wrench y límites};
\node[bx,fill=softwarm](r)at(65,10){\textbf{Rigid attachment}\\formación rígida\\distribución de fuerza};
\node[bx,fill=softgreen](p)at(109,10){\textbf{Pushing / caging}\\contacto unilateral\\geometría no-prehensil};
\node[bx,fill=softplum](g)at(153,10){\textbf{Grasping / arms}\\grasp matrix\\force closure};
\foreach \a/\b in {s/r,r/p,p/g}{\draw[-{Stealth[length=1.55mm]},muted,line width=.42pt](\a)--(\b);}
\end{tikzpicture}%
  }}
  \caption[Matriz de capacidades y modalidades físicas.]{(a) Matriz de
  capacidades de los trabajos discriminantes: \capyes\ evidencia explícita,
  \cappart\ cobertura parcial o condicionada y \capno\ no documentado en el
  alcance revisado. La columna Ev. indica el tipo de evidencia (T teórica,
  S simulación, P plataforma física) y la columna C el nivel de centralización,
  de C0 ---decisión centralizada--- a C4 ---decisión local con intercambios
  vecinales sin autoridad central en tiempo de ejecución---. (b) Modalidades
  físicas que no deben mezclarse bajo la etiqueta genérica \emph{cooperative
  transport}: cada una impone un conjunto de esfuerzos admisibles distinto.}
  \label{fig:lit-capability-matrix}
  \viusource{elaboración propia a partir de la codificación del núcleo canónico}
\end{figure}


La Figura~\ref{fig:lit-capability-matrix} se lee por el corte horizontal, no por
columnas aisladas: las seis primeras filas cubren decisión de coalición con
arquitecturas distribuidas y ninguna toca la carga compartida; las seis siguientes
cubren carga compartida con certificación física y ninguna recluta desde una flota
mayor. La columna de centralización añade el matiz decisivo, porque separa los trabajos
que deciden localmente de los que conservan una autoridad central en tiempo de
ejecución. Un círculo vacío señala que la fuente no documenta esa capacidad, no que el
método sea incapaz de incorporarla.

Tres lecturas condicionan el diseño del resto de la memoria. La atomicidad exige que el
cierre entero pertenezca al método, porque los juegos poblacionales producen masa
continua y un redondeo posterior puede romper la propiedad demostrada en la
Ecuación~\eqref{eq:tf-population-dynamics}. El paso de capacidad a \emph{wrench} impide
sustituir $W_k^{\mathrm{req}} \in \mathcal W_C(q)$ por la suma escalar
$\sum_i c_i \ge d_k$, según la Ecuación~\eqref{eq:tf-grasp-map}. La escala obliga a
separar robots simulados, robots físicos, robots por coalición y cargas simultáneas.
A esto se añade la modalidad de interacción: como recoge el panel inferior de la
Figura~\ref{fig:lit-capability-matrix}, carga soportada, acoplamiento rígido, empuje y
prensión imponen conjuntos de esfuerzos admisibles distintos y no deben agruparse.

% Tabla 2 del consolidado académico (PDF vigente, página 4): antecedentes más
% próximos frente a las ocho capacidades que definen la combinación objetivo.
% Reconstruida en LaTeX; el `.tex` depositado corresponde a una versión previa
% con otra tabla de red-team.
\begin{table}[!htbp]
  \centering
  \caption{Antecedentes más próximos frente a las capacidades que definen la combinación objetivo.}
  \label{tab:lit-prior-art}
{\setstretch{1}\fontsize{7.6}{8.6}\selectfont\setlength{\tabcolsep}{3.4pt}\renewcommand{\arraystretch}{1.22}
\rowcolors{2}{white}{soft}
\begin{tabularx}{\linewidth}{@{}J{2.6cm}*{8}{K{0.62cm}}Z@{}}
\toprule
\textbf{Trabajo} & \textbf{H} & \textbf{V} & \textbf{L} & \textbf{S} & \textbf{W} & \textbf{D} & \textbf{R} & \textbf{X} & \textbf{Restricción decisiva}\\
\midrule
Verma 2019 \parencite*{verma2019replacement} & \capno & \cappart & \capno & \capyes & \cappart & \capno & \capyes & \capyes & planificación global; ruta e hitos conocidos\\
Sahu 2026 \parencite*{sahuKumar2026FaultManagement} & \capno & \cappart & \capno & \capyes & \cappart & \cappart & \capyes & \capyes & arquitectura maestro--esclavo\\
Shibata 2023 \parencite*{shibata2023event} & \capno & \cappart & \capno & \capyes & \capyes & \capyes & \cappart & \cappart & equipo variable; no selección desde una flota heterogénea\\
Ebel 2024 \parencite*{ebel2024cooperative} & \capno & \capno & \capno & \capyes & \capyes & \capyes & \capno & \capno & equipo de transporte fijado\\
Bezerra 2025 \parencite*{bezerra2025dynamicCoalition} & \capyes & \capyes & \capyes & \capno & \capno & \capyes & \cappart & \capno & reclutamiento lógico; sin carga compartida física\\
GRAPE-S 2026 \parencite*{diehlAdams2026Grapes} & \capyes & \capyes & \cappart & \capno & \capno & \capyes & \capno & \capno & formación de servicios; no transporte de carga compartida\\
\bottomrule
\end{tabularx}}
\par\vspace{1.5pt}
{\setstretch{1}\fontsize{7.0}{7.8}\selectfont\color{muted}
H = heterogeneidad; V = tamaño de coalición variable; L = reclutamiento local; S = carga compartida; W = certificado físico; D = ejecución o decisión distribuida; R = recuperación; X = sustitución de miembro. \capyes\ explícito, \cappart\ parcial o condicionado, \capno\ no documentado en el alcance revisado.\par}
  \viusource{codificación del núcleo canónico y evidencia de texto completo}
\end{table}


Una brecha solo es defendible si sobrevive al intento de refutarla, y la
Tabla~\ref{tab:lit-prior-art} recoge esa auditoría adversarial. El resultado obliga a
estrechar el enunciado: el reemplazo durante el transporte ya existe. Verma et al. lo
implementan con ruta y puntos de relevo conocidos y planificación global
\parencite{verma2019replacement}; Sahu y Kumar lo integran con detección de fallo bajo
arquitectura maestro--esclavo \parencite{sahuKumar2026FaultManagement}; y Shibata et al.
admiten que agentes se retiren o incorporen durante el control distribuido, sin resolver
la selección del sustituto dentro de una flota heterogénea \parencite{shibata2023event}.
En el extremo opuesto, Bezerra y GRAPE-S cubren reclutamiento o formación de coaliciones
sin cerrar la interacción física de una carga compartida.

\begin{tcolorbox}[enhanced,breakable,colback=VIUOrange!5,colframe=VIUOrange,
  boxrule=0.45pt,leftrule=2.8pt,arc=1.2mm,left=2.5mm,right=2.5mm,top=1.8mm,bottom=1.8mm]
\textbf{Brecha académica defendible.} En el corpus revisado no se identificó una
arquitectura que integre, en una misma cadena operativa, el reclutamiento local de una
coalición heterogénea de tamaño variable desde una flota mayor, la certificación
mecánica de la carga compartida y el reemplazo en línea de un miembro durante el
transporte, manteniendo la decisión de coalición y la ejecución sin una autoridad global
permanente. El enunciado se restringe al protocolo de búsqueda documentado y no afirma
inexistencia absoluta.
\end{tcolorbox}

El valor del enunciado está en lo que excluye: no reclama el primer método distribuido
de transporte cooperativo, ni el primer mecanismo de reemplazo, ni superioridad frente a
ningún trabajo de la Figura~\ref{fig:lit-capability-matrix}. Reclama solo que la
composición de las tres piezas bajo información local carece de precedente documentado
en el corpus, y esa es la afirmación que el capítulo de resultados evalúa por tramos.
